% Draft for the conference paper (full paper deadline: 10 October).
% Fill in authors/affiliations, check the conference template, then paste
% the figures from ../results.
\documentclass[conference]{IEEEtran}
\usepackage{amsmath,amssymb,graphicx,booktabs}
\begin{document}

\title{Adaptive Passivity-Based Control with Online ANN Gain Tuning of a
2~MW Direct-Drive PMSG Wind Energy Conversion System}

\author{\IEEEauthorblockN{Said Regani, [Supervisor]}
\IEEEauthorblockA{[Laboratory], [University], Algeria}}

\maketitle

\begin{abstract}
This paper proposes a hybrid control strategy for a grid-connected 2~MW
direct-drive permanent-magnet synchronous generator (PMSG) wind turbine
equipped with a back-to-back converter and an L filter. Passivity-based
control (PBC), which exploits the workless (skew-symmetric) coupling of the
machine and filter models, is combined with Lyapunov-based adaptation of the
uncertain electrical parameters (stator resistance and inductance, flux
linkage, filter resistance and inductance), of the aerodynamic torque and of
the DC-link power disturbance. The damping-injection gains of every loop are
then tuned online by radial-basis-function neural networks. Because the
closed-loop storage function decreases for any positive damping gain, the
neural tuner improves transients without affecting the stability proof.
Simulations under wind variations, parameter mismatch and a 20\% grid voltage
dip show [X]\% lower DC-link voltage deviation than conventional PI vector
control, with identical MPPT performance.
\end{abstract}

\begin{IEEEkeywords}
PMSG, wind energy, back-to-back converter, passivity-based control, adaptive
control, neural network, MPPT, DC-link voltage.
\end{IEEEkeywords}

\section{Introduction}
% - PMSG direct drive, back-to-back converter: context, advantages.
% - PI vector control: sensitive to parameter variations, tuning.
% - PBC (Ortega et al.) used for PMSG / grid converters; adaptive control used
%   separately; their combination with online ANN tuning has not been reported.
% - Contributions: (i) unified adaptive PBC for MSC, DC bus and GSC;
%   (ii) Lyapunov proof valid for time-varying gains; (iii) RBF online tuning;
%   (iv) comparison with PI on a 2 MW system.

\section{Modelling of the Wind Energy Conversion System}
\subsection{Turbine}
\begin{equation}
P_w=\tfrac12\rho\pi R^2 C_p(\lambda,\beta)v^3,\quad \lambda=\frac{\Omega R}{v},\quad
J\dot\Omega=T_w-T_e-B\Omega
\end{equation}
\subsection{PMSG}
With $\mathbf{J}=\left[\begin{smallmatrix}0&1\\-1&0\end{smallmatrix}\right]$ and
$\mathbf e_q=[0\;1]^T$:
\begin{equation}
L_s\dot{\mathbf i}_s=-R_s\mathbf i_s+\omega_eL_s\mathbf J\mathbf i_s+\omega_e\psi_f\mathbf e_q-\mathbf v_s,\quad
T_e=\tfrac32 p\psi_f i_{sq}
\end{equation}
\subsection{DC link and grid filter}
\begin{equation}
\dot W=P_{msc}-P_{gsc},\; W=\tfrac12CV_{dc}^2;\quad
L_f\dot{\mathbf i}_g=-R_f\mathbf i_g+\omega_gL_f\mathbf J\mathbf i_g+\mathbf v_i-\mathbf v_g
\end{equation}

\section{Proposed Adaptive PBC with ANN Gain Tuning}
\subsection{Machine-side current loop}
\begin{equation}
\mathbf v_s=-\hat L_s\dot{\mathbf i}_s^*-\hat R_s\mathbf i_s^*+\omega_e\hat L_s\mathbf J\mathbf i_s^*
+\omega_e\hat\psi\mathbf e_q+K_a(t)\mathbf e_s
\end{equation}
\begin{equation}
\dot{\hat R}_s=-\gamma_R\mathbf e_s^T\mathbf i_s^*,\;
\dot{\hat L}_s=\gamma_L\mathbf e_s^T(\omega_e\mathbf J\mathbf i_s^*-\dot{\mathbf i}_s^*),\;
\dot{\hat\psi}=\gamma_\psi\omega_e e_{sq}
\end{equation}
\textbf{Proposition.} With
$V=\tfrac12\mathbf e_s^TL_s\mathbf e_s+\sum\tilde\theta^2/(2\gamma_\theta)$,
$\dot V=-(R_s+K_a(t))\|\mathbf e_s\|^2\le0$ for any $K_a(t)>0$.
% proof: substitute, use e^T J e = 0 (see docs/modelisation_commande.md)
\subsection{Speed (MPPT) loop}
$T_e^*=\hat T_w-B\Omega-J\dot\Omega^*+k_\Omega(t)e_\Omega$, $\dot{\hat T}_w=\gamma_Te_\Omega$,
$\Omega^*=\lambda_{opt}v/R$.
\subsection{DC-link energy loop and grid-side current loop}
$P_{gsc}^*=P_{msc}+\hat d+k_v(t)e_W$, $\dot{\hat d}=\gamma_de_W$;
$\mathbf v_i=\hat L_f\dot{\mathbf i}_g^*+\hat R_f\mathbf i_g^*-\omega_g\hat L_f\mathbf J\mathbf i_g^*+\mathbf v_g-K_b(t)\mathbf e_g$.
\subsection{Online RBF tuning of the damping gains}
$K(t)=g(t)K_0$, $g=g_{min}+(g_{max}-g_{min})\,\mathrm{sig}(\mathbf w^T\boldsymbol\phi(e_n,\dot e_n))$,
$\dot{\mathbf w}=\eta e_n^2\,\partial g/\partial\mathbf w-\sigma\mathbf w$.

\section{Simulation Results}
% Table I: system parameters (docs/modelisation_commande.md, Sec. 1)
% Fig. 1: system + control block diagram
% Fig. 2: results/fig_nominal.png ; Fig. 3: results/fig_robust.png
% Fig. 4: results/fig_robust_estimates.png
% Table II: indices printed by matlab/main.m

\section{Conclusion}

\bibliographystyle{IEEEtran}
% \bibliography{refs}
\end{document}
